
\paragraph{\textbf{Corollary 6.2.}}
\SGARefTarget{sga3:xiv:6.2}
{\itshape Let \(G\) be a smooth solvable algebraic group over a
field \(k\). Then the variety \(\mathcal C\) of Cartan subgroups
of \(G\) is isomorphic to an affine space
\(\operatorname{Spec}k[t_1,\ldots,t_N]\).}

We may suppose that \(G\) is connected and affine. Let
\(G_\infty\) be the terminal group in the descending central
series
\[
  G_0=G,\qquad G_{i+1}=[G,G_i].
\]
Thus \(G_\infty\) is the smallest invariant algebraic subgroup
of \(G\) for which \(G/G_\infty\) is nilpotent. Let \(C\) be a
Cartan subgroup of \(G\), which exists by \SGARefLink{sga3:xiv:1.1}{1.1}. Its image in
\(G/G_\infty\) is a Cartan subgroup and hence is all of
\(G/G_\infty\). Consequently the morphism
\[
  G_\infty\longrightarrow\mathcal C,\qquad
  g\longmapsto\operatorname{Ad}(g)C
\]
is an epimorphism and identifies
\[
  \mathcal C\simeq G_\infty/(N\cap G_\infty),
\]
where \(N\) is the normalizer of \(C\) in \(G\). In fact \(N=C\),
as recalled in \SGARefLink{sga3:xiv:4.4}{4.4}, but that is immaterial here.

The group \(U=G_\infty\) is smooth, connected, and unipotent:
over \(\bar k\) it lies in the unipotent part of \(G\), by the
known structure of smooth affine solvable groups
\textup{\emph{Bible}~6, Theorem~3}. If \(k\) is perfect, as in
the finite-field case needed for \SGARefLink{sga3:xiv:6.1}{6.1}, it follows easily that
\(U\) is \(k\)-unipotent: it has a composition series by
algebraic subgroups \(U_i\) with
\(U_i/U_{i+1}\simeq\mathbb G_a\).

Rosenlicht proved that a group of the form \(G_\infty\) above is
\(k\)-unipotent without any restriction on \(k\)
\textup{(M.~Rosenlicht, \emph{Questions of rationality for
solvable algebraic groups over non perfect fields}, Annali di
Matematica, 1963, pp.~97--120, \SGARefLink{sga3:xiv:7:lemma4}{Theorem~4}, \SGARefLink{sga3:xiv:7:lemma2}{Corollary~2})}. We admit
this substantially more delicate result and apply the following
lemma.

\paragraph{\textbf{Lemma 6.3.}}
\SGARefTarget{sga3:xiv:6.3}
{\itshape Let \(U\) be a smooth connected algebraic group over a
field \(k\), and let \(X=U/V\) be a homogeneous space under \(U\)
having a \(k\)-rational point. If \(U\) is \(k\)-unipotent, then,
as a \(k\)-scheme, \(X\) is isomorphic to an affine space
\(\operatorname{Spec}k[t_1,\ldots,t_N]\).}

Choose a composition series
\[
  U=U_0\supset U_1\supset\cdots\supset U_n=(e)
\]
by smooth connected normal subgroups of \(U\), with
\(U_i/U_{i+1}\simeq\mathbb G_a\). The groups
\[
  K_i=U_iV
\]
are algebraic subgroups of \(U\). They need not be smooth or
connected when \(V\) is not, but this does not matter, and
\(K_{i+1}\) is normal in \(K_i\). The canonical epimorphism
\[
  U_i/U_{i+1}\longrightarrow K_i/K_{i+1}
\]
shows that \(K_i/K_{i+1}\) is either the unit group or isomorphic
to \(\mathbb G_a/H_i\), where \(H_i\) is a finite subgroup of
\(\mathbb G_a\). Rosenlicht's \SGARefLink{sga3:xiv:7:lemma2}{Theorem~2} in the cited paper gives
\[
  K_i/K_{i+1}\simeq\mathbb G_a;
\]
this last assertion is immediate when \(k\) is perfect.

Put \(X_i=U/K_i\). We prove inductively that every \(X_i\) is an
affine space. If \(K_i/K_{i+1}\) is trivial, then
\(X_i=X_{i+1}\). Otherwise \(X_{i+1}\) is a principal
\(\mathbb G_a\)-bundle over \(X_i\). Since \(X_i\) is affine,
the bundle is trivial, so
\[
  X_{i+1}\simeq X_i\times\mathbb G_a.
\]
The induction proves the lemma. It proves \SGARefLink{sga3:xiv:6.2}{6.2} and thereby
completes the proof of \SGARefLink{sga3:xiv:6.1}{6.1}.
